\subsubsection{Quintic order, \texorpdfstring{$O(\gs^3)$}{O(gs**3)}}
\label{p:sec:hs:quintic}

\paragraph{Quintic coupling and equations of motion}

\subparagraph{The quintic coupling.}

The next term of the action is the quintic term:
\begin{equation}
  \label{p:eq:hs5:S5}
  S_5
    = \frac{\gs^3}{5!} \, \vertex{5}(\Psi^5),
\end{equation}
where the vertex $\vertex{5}$ has five slots.
The question is how to describe this vertex in terms of auxiliary fields.
We cannot use the same trick as for the quartic order, because it requires an exact square.

Given the form of the saddle-point solution \eqref{p:eq:hs4:saddle} for $\Sigma_2$, we can replace $\Psi^4$ by $\Sigma_2^2$ to turn the quintic term into a cubic coupling.
Two auxiliary fields $\Sigma_2$ fill four of the five slots and the string field fills the last one:
\begin{equation}
  \label{p:eq:hs5:S5tilde}
  \caction[5]
    := \frac{\gs^3}{5!} \, \vertex{5}(\Sigma_2^2, \Psi),
  \qquad
  \caction[\le 5]
    := \caction[\le 4] + \caction[5],
\end{equation}
where $\vertex{5}(\Sigma_2^2, \Psi) := \vertex{5}(\Sigma_2, \Sigma_2, \Psi)$.
At this stage this should be treated as an ansatz: we will now check that eliminating $\Sigma_2$ gives back $S_{\le 5}$.
Note that the partition $5 = 2 + 2 + 1$ of the five strands over three fields drawn from $\{\Sigma_2, \Psi\}$ is the only one, so no other cubic term can be built from the quintic vertex with these two fields.
In principle we could use a $3$-string field to partition $5 = 3 + 1 + 1$, but it is not needed before the sextic order: we want to keep the field content minimal and avoid introducing fields from higher orders into lower orders.
Such a term also appeared when decomposing the quintic term of open SFT with stubs~\cite{Erbin:2024:OpenStringStub}.

Ghost number is conserved as it must: the two $\Sigma_2$ and the $\Psi$ carry $4 + 4 + 2 = 10$, the same total as $\Psi^5$.
Since the vertex is read with two of its slots as outputs, we define the product $\sprod{3}{2} \colon \Hn{3} \to \Hn{2}$ by
\begin{equation}
  \label{p:eq:hs5:p32}
  \vertex{5}(A_1, A_2, A_3, A_4, A_5)
    =: \pair[2]{A_1 \otimes A_2}{\sprod{3}{2}(A_3, A_4, A_5)},
\end{equation}
as $\sprod{2}{2}$ was defined from the quartic vertex, with the $2$-string pairing on the left.
It has ghost number zero, by the general formula~\eqref{p:eq:conv:ghost-numbers}.

Since the cubic term $\caction[5]$ contains $\Sigma_2$, it modifies the equation of motion of $\Sigma_2$, whose solution is no longer $\Psi^2$.
The terms of $\caction[\le 5]$ which contain $\Sigma_2$ are:
\begin{equation}
  \label{p:eq:hs5:F}
  \caction[4] + \caction[5]
    = - \frac{\gs^2}{4!} \, \vertex{4}(\Sigma_2^2)
      + \frac{2 \gs^2}{4!} \, \vertex{4}(\Sigma_2, \Psi^2)
      + \frac{\gs^3}{5!} \, \vertex{5}(\Sigma_2^2, \Psi),
\end{equation}
a quadratic form in $\Sigma_2$ with two terms, and a source term linear in $\Sigma_2$.
The kinetic term and the cubic vertex of $\Psi$ are the same in $\caction[\le 5]$ and in $S_{\le 5}$ and drop out of every comparison.


\subparagraph{Equations of motion.}

Varying the cubic action~\eqref{p:eq:hs5:S5tilde} and writing each term through the pairing, we get:
\begin{equation}
  \label{p:eq:hs5:delta-S}
  \begin{aligned}
    \delta \caction[\le 5]
      &
      = \pair{\delta \Psi}{Q_B \Psi}
        + \frac{\gs}{2} \, \pair{\delta \Psi}{\ell_2(\Psi^2)}
        + \frac{\gs^2}{3!} \, \pair{\delta \Psi}{\ell_3(\Psi, \Sigma_2)}
        + \frac{\gs^3}{5!} \, \pair{\delta \Psi}{\ell_4(\Sigma_2^2)}
        \\ &
        \quad
        - \frac{2 \gs^2}{4!} \, \pair[2]{\delta \Sigma_2}{\sprod{2}{2}(\Sigma_2)}
        + \frac{2 \gs^2}{4!} \, \pair[2]{\delta \Sigma_2}{\sprod{2}{2}(\Psi^2)}
        + \frac{2 \gs^3}{5!} \, \pair[2]{\delta \Sigma_2}{\sprod{3}{2}(\Psi, \Sigma_2)}.
  \end{aligned}
\end{equation}
The two new terms are the variation of~\eqref{p:eq:hs5:S5tilde} with respect to $\Psi$ and to $\Sigma_2$.
The equations of motion are then:
\begin{subequations}
  \label{p:eq:hs5:eom}
  \begin{align}
    \label{p:eq:hs5:eom-psi}
    \mathcal E^{(3)}_\Psi
      &
      := Q_B \Psi
        + \frac{\gs}{2} \, \ell_2(\Psi^2)
        + \frac{\gs^2}{3!} \, \ell_3(\Psi, \Sigma_2)
        + \frac{\gs^3}{5!} \, \ell_4(\Sigma_2^2)
      = 0,
    \\
    \label{p:eq:hs5:eom-sigma}
    \mathcal E^{(3)}_{\Sigma_2}
      &
      := \sprod{2}{2}(\Psi^2)
        - \sprod{2}{2}(\Sigma_2)
        + \frac{\gs}{5} \, \sprod{3}{2}(\Psi, \Sigma_2)
      = 0.
  \end{align}
\end{subequations}
The equation for $\Sigma_2$ is still algebraic, but it no longer states that $\Sigma_2$ and $\Psi^2$ have the same image under $\sprod{2}{2}$: the quintic vertex adds a term linear in $\Sigma_2$ and in $\Psi$.
We group this term with the quadratic form, since both are linear in $\Sigma_2$, and define the operator:
\begin{equation}
  \label{p:eq:hs5:dressed}
  \kin{2}{\Psi}{3}
    := \sprod{2}{2} - \frac{\gs}{5} \, \sprod{3}{2}(\Psi, \cdot),
  \qquad
  \kin{2}{\Psi}{3} \colon \Hn{2} \to \Hn{2},
\end{equation}
which we call the \emph{dressed kinetic operator} of $\Sigma_2$.
The quintic vertex, with one slot filled by $\Psi$, is not a kinetic term: it dresses the kinetic operator $\sprod{2}{2}$, as interaction insertions dress a propagator.
The superscript records that it is the operator of the action truncated at order $\gs^3$, since it receives further corrections at the next orders.
The dressed operator is symmetric with respect to the $2$-string pairing, $\pair[2]{A}{\kin{2}{\Psi}{3}(B)} = \pair[2]{\kin{2}{\Psi}{3}(A)}{B}$, since both vertices are symmetric, and it is affine in $\Psi$.
In terms of it, the quadratic part of~\eqref{p:eq:hs5:F} is a single form:
\begin{equation}
  \label{p:eq:hs5:F-dressed}
  - \frac{\gs^2}{4!} \, \vertex{4}(\Sigma_2^2)
    + \frac{\gs^3}{5!} \, \vertex{5}(\Sigma_2^2, \Psi)
    = - \frac{\gs^2}{4!} \, \pair[2]{\Sigma_2}{\kin{2}{\Psi}{3}(\Sigma_2)}.
\end{equation}
The equation of motion of $\Sigma_2$ then reads:
\begin{equation}
  \label{p:eq:hs5:eom-sigma-dressed}
  \mathcal E^{(3)}_{\Sigma_2}
    = \sprod{2}{2}(\Psi^2) - \kin{2}{\Psi}{3}(\Sigma_2)
    = 0.
\end{equation}
We can make an analogy with non-linear sigma models: the quadratic form of $\Sigma_2$ has a $\Psi$-dependent metric, and the field $\Psi$ enters the kinetic term of $\Sigma_2$ as a background.


\paragraph{Saddle-point solution and induced interactions}

\subparagraph{Saddle-point solution.}

On the quotient, where $\sprod{2}{2}$ is injective, the dressed operator is $\sprod{2}{2}$ composed with a perturbation of the identity:
\begin{equation}
  \label{p:eq:hs5:dressed-factor}
  \kin{2}{\Psi}{3}
    = \sprod{2}{2} \circ \Big( 1 - \frac{\gs}{5} \, \sprod{2}{2}^{-1} \circ \sprod{3}{2}(\Psi, \cdot) \Big),
\end{equation}
provided that $\sprod{3}{2}(\Psi, \cdot)$ takes its values in the image of $\sprod{2}{2}$ and descends to the quotient.
This is an assumption on the quintic vertex, stated precisely in~\eqref{p:eq:hs5:kernel-condition} below, where the two standing assumptions of the construction from this order on are collected.
For small $\gs$, the inverse of the dressed operator then exists as a Neumann series:
\begin{equation}
  \label{p:eq:hs5:neumann}
  \big( \kin{2}{\Psi}{3} \big)^{-1}
    = \sum_{n \ge 0} \Big( \frac{\gs}{5} \, \sprod{2}{2}^{-1} \circ \sprod{3}{2}(\Psi, \cdot) \Big)^n \circ \sprod{2}{2}^{-1},
\end{equation}
which converges for $\gs$ small at fixed $\Psi$ when the composite $\sprod{2}{2}^{-1} \circ \sprod{3}{2}(\Psi, \cdot)$ is bounded on the quotient.
Order by order in $\gs$, no invertibility condition appears beyond the one on $\sprod{2}{2}$ assumed in \cref{p:sec:hs:quartic}.
The statements at finite $\gs$ below assume instead that $\kin{2}{\Psi}{3}$ itself is invertible on the quotient.
The equation of motion~\eqref{p:eq:hs5:eom-sigma-dressed} then has the unique solution:
\begin{equation}
  \label{p:eq:hs5:saddle}
  \Sigma_2^{(3)*}
    := \big( \kin{2}{\Psi}{3} \big)^{-1} \circ \sprod{2}{2}(\Psi^2)
    = \Psi^2 + \frac{\gs}{5} \, \big( \kin{2}{\Psi}{3} \big)^{-1} \circ \sprod{3}{2}(\Psi^3),
\end{equation}
where the second form follows from $\sprod{2}{2}(\Psi^2) = \kin{2}{\Psi}{3}(\Psi^2) + \tfrac{\gs}{5} \, \sprod{3}{2}(\Psi, \Psi^2)$, the definition~\eqref{p:eq:hs5:dressed} applied to $\Psi^2$.
The solution is again fixed entirely by $\Psi$ through an algebraic equation, so that $\Sigma_2$ remains an auxiliary field.
It is no longer $\Psi^2$ but $\Psi^2$ plus corrections at all orders in $\gs$.
The solution now contains more terms, and we could ask whether the correction of order $\gs$ spoils the matching, but this is not the case by construction.
By stationarity of the quartic-order action at $\Psi^2$, the action is quadratic in any fluctuation around it, and the quintic term, of order $\gs^3$, changes linearly in the correction of order $\gs$, so the correction to the eliminated action starts at order $\gs^4$, as we will see below.

Applying $\sprod{2}{2}^{-1}$ to~\eqref{p:eq:hs5:eom-sigma} instead gives the solution as a recursion:
\begin{equation}
  \label{p:eq:hs5:recursion}
  \begin{aligned}
    \Sigma_2^{(3)*}
      &
      = \Psi^2 + \frac{\gs}{5} \, \sprod{2}{2}^{-1} \circ \sprod{3}{2}(\Psi, \Sigma_2^{(3)*})
      \\ &
      = \quad
        \begin{tikzpicture}[diagram]
          \draw (0, 0.1) -- (2, 0.1) pic {ext} node[above right] {$\Psi$};
          \draw (0, -0.1) -- (2, -0.1) pic {ext} node[below right] {$\Psi$};
        \end{tikzpicture}
        \quad
        +
        \quad
        \begin{tikzpicture}[diagram]
          \draw[sigma line] (0, 0.1) -- (4, 0.1) pic {ext};
          \draw[sigma line] (0, -0.1) -- (4, -0.1) pic {ext};
          \node[right=3pt] at (4, 0) {$\Sigma_2^{(3)}$};
          \draw (2, 0.1) -- (2, 1.4) pic {ext} node[above=2pt] {$\Psi$};
          \node[below=3pt] at (0.9, -0.1) {$\sprod{2}{2}^{-1}$};
          \node[below=3pt] at (2, -0.1) {$\sprod{3}{2}$};
        \end{tikzpicture}
  \end{aligned}
\end{equation}
Iterating it, we get:
\begin{equation}
  \label{p:eq:hs5:saddle-series}
  \begin{aligned}
    \Sigma_2^{(3)*}
      &
      = \Psi^2
        + \frac{\gs}{5} \, \sprod{2}{2}^{-1} \circ \sprod{3}{2}(\Psi^3)
        + \frac{\gs^2}{25} \, \sprod{2}{2}^{-1} \circ \sprod{3}{2}\Big( \Psi, \sprod{2}{2}^{-1} \circ \sprod{3}{2}(\Psi^3) \Big)
        + O(\gs^3)
      \\ &
      = \quad
        \begin{tikzpicture}[diagram]
          \draw (0, 0.1) -- (2, 0.1) pic {ext} node[above right] {$\Psi$};
          \draw (0, -0.1) -- (2, -0.1) pic {ext} node[below right] {$\Psi$};
        \end{tikzpicture}
        \quad
        +
        \quad
        \begin{tikzpicture}[diagram]
          \draw[sigma line] (0, 0.1) -- (2, 0.1) pic {ext};
          \draw[sigma line] (0, -0.1) -- (2, -0.1) pic {ext};
          \draw (2, 0.1) -- (4, 0.1) pic {ext} node[above right] {$\Psi$};
          \draw (2, -0.1) -- (4, -0.1) pic {ext} node[below right] {$\Psi$};
          \draw (2, 0.1) -- (2, 1.4) pic {ext} node[above=2pt] {$\Psi$};
          \node[below=3pt] at (0.9, -0.1) {$\sprod{2}{2}^{-1}$};
          \node[below=3pt] at (2, -0.1) {$\sprod{3}{2}$};
        \end{tikzpicture}
        \quad
        +
        \quad
        \cdots
  \end{aligned}
\end{equation}
which is the Neumann series~\eqref{p:eq:hs5:neumann} applied to $\sprod{3}{2}(\Psi^3)$ (we did not include the numerical factors in the graphical representations).
Each iteration inserts one more $\sprod{3}{2}(\Psi, \cdot)$ between two $\sprod{2}{2}^{-1}$, with one more factor $\gs / 5$.
Read as Feynman diagrams, the term of order $\gs^n$ is a chain of $n$ quintic vertices joined by $\Sigma_2$ lines, each carrying one external $\Psi$.
The chain ends on the source $\Psi^2$ at one end and on the open $\Sigma_2$ slot at the other.
Up to the factor $12 / \gs^2$, the inverse of the dressed operator is the $\Sigma_2$ propagator in the background $\Psi$, and the chain is the geometric series of the insertions $\sprod{3}{2}(\Psi, \cdot)$ on the $\Sigma_2$ line.

The field $\Sigma_2$ enters the cubic action through its fluctuation from the solution:
\begin{equation}
  \label{p:eq:hs5:deviation}
  \sigma_2^{(3)}
    := \Sigma_2 - \Sigma_2^{(3)*},
\end{equation}
which replaces the fluctuation $\sigma_2^{(2)} = \Sigma_2 - \Psi^2$ of the quartic order.
The terms~\eqref{p:eq:hs5:F} can be written as a completed square in the fluctuation.
By~\eqref{p:eq:hs5:F-dressed} and the definition~\eqref{p:eq:hs5:saddle} of the solution, $\sprod{2}{2}(\Psi^2) = \kin{2}{\Psi}{3}(\Sigma_2^{(3)*})$, we get:
\begin{equation}
\label{p:eq:hs5:square}
\begin{aligned}
  \caction[4] + \caction[5]
    &
    = - \frac{\gs^2}{4!} \Big[ \pair[2]{\Sigma_2}{\kin{2}{\Psi}{3}(\Sigma_2)} - 2 \, \pair[2]{\Sigma_2}{\kin{2}{\Psi}{3}(\Sigma_2^{(3)*})} \Big]
    \\ &
    = - \frac{\gs^2}{4!} \Big[ \pair[2]{\sigma_2^{(3)}}{\kin{2}{\Psi}{3}(\sigma_2^{(3)})} - \pair[2]{\Sigma_2^{(3)*}}{\kin{2}{\Psi}{3}(\Sigma_2^{(3)*})} \Big],
\end{aligned}
\end{equation}
using the symmetry of the dressed operator.
The first term is the quadratic form of the fluctuation and vanishes on-shell.
The second term does not depend on $\Sigma_2$ and is the value of the action on the solution.


\subparagraph{Induced interactions.}

Expanding the second form of~\eqref{p:eq:hs5:saddle} in the constant term of~\eqref{p:eq:hs5:square}, and using $\kin{2}{\Psi}{3} \circ (\kin{2}{\Psi}{3})^{-1} = 1$ together with $\pair[2]{A}{\sprod{3}{2}(\Psi^3)} = \vertex{5}(A, \Psi^2, \Psi)$, we get:
\begin{equation}
  \label{p:eq:hs5:constant}
  \pair[2]{\Sigma_2^{(3)*}}{\kin{2}{\Psi}{3}(\Sigma_2^{(3)*})}
    = \vertex{4}(\Psi^4)
      + \frac{\gs}{5} \, \vertex{5}(\Psi^5)
      + \frac{\gs^2}{25} \, \vertex{5}\Big( \Psi^3, \big( \kin{2}{\Psi}{3} \big)^{-1} \circ \sprod{3}{2}(\Psi^3) \Big).
\end{equation}
The first two terms come from $\pair[2]{\Psi^2}{\kin{2}{\Psi}{3}(\Psi^2)} = \vertex{4}(\Psi^4) - \tfrac{\gs}{5} \vertex{5}(\Psi^5)$ and from twice the cross term $\tfrac{\gs}{5} \vertex{5}(\Psi^5)$.
Multiplied by $\gs^2 / 4!$, the first two terms are the quartic and quintic terms of $S_{\le 5}$.
This is how the coefficient $1 / 5!$ of~\eqref{p:eq:hs5:S5tilde} enters: the quintic term of the target is the cross term minus the quintic part of the dressed operator, and the two combine to $+ \gs / 5$ only with this coefficient.
Putting~\eqref{p:eq:hs5:square} and~\eqref{p:eq:hs5:constant} together, and adding back the terms of $\Psi$ alone, we get:
\begin{equation}
  \label{p:eq:hs5:identity}
  \caction[\le 5]
    = S_{\le 5}
      - \frac{\gs^2}{4!} \, \pair[2]{\sigma_2^{(3)}}{\kin{2}{\Psi}{3}(\sigma_2^{(3)})}
      + \frac{\gs^4}{600} \, \vertex{5}\Big( \Psi^2, \Psi, \big( \kin{2}{\Psi}{3} \big)^{-1} \circ \sprod{3}{2}(\Psi^3) \Big),
\end{equation}
This is an identity in the fields $(\Psi, \Sigma_2)$, exact in $\gs$, and the counterpart at this order of the completed square~\eqref{p:eq:hs4:action-tilde}.
The three terms are the target, a quadratic form in the fluctuation which vanishes on-shell, and a term independent of $\Sigma_2$.
The last term starts at order $\gs^4$, where it is of degree six in $\Psi$.
It is a six-point interaction induced by the exchange of $\Sigma_2$, which the Gaussian elimination produces in addition to the vertex for which the auxiliary field was introduced.

On-shell, $\sigma_2^{(3)} = 0$ and the elimination of $\Sigma_2$ gives:
\begin{equation}
  \label{p:eq:hs5:eliminate}
  \caction[\le 5] \big|_{\Sigma_2 = \Sigma_2^{(3)*}}
    = S_{\le 5}
      + \frac{\gs^4}{600} \, \vertex{5 \cdot 5}(\Psi^6)
      + O(\gs^5),
\end{equation}
where we expanded the dressed inverse with~\eqref{p:eq:hs5:neumann} and introduced the composite vertex
\begin{equation}
  \label{p:eq:hs5:exchange-vertex}
  \begin{aligned}
    \vertex{5 \cdot 5}(A_1, \ldots, A_6)
      &
      := \vertex{5}\Big( A_1, A_2, A_3, \sprod{2}{2}^{-1} \circ \sprod{3}{2}(A_4, A_5, A_6) \Big)
      \\ &
      = \pair[2]{\sprod{3}{2}(A_1, A_2, A_3)}{\sprod{2}{2}^{-1} \circ \sprod{3}{2}(A_4, A_5, A_6)},
  \end{aligned}
\end{equation}
a composite six-point vertex obtained from two quintic vertices joined by a $\Sigma_2$ line.
It is a six-linear form, graded-symmetric within each triple of arguments and under the exchange of the two triples, since $\sprod{2}{2}^{-1}$ is symmetric with respect to the $2$-string pairing as $\sprod{2}{2}$ is.
We read it also as a trilinear form on $2$-string states: $\vertex{5 \cdot 5}(B_1, B_2, B_3)$ with $B_i \in \Hn{2}$ places $B_1$ in the first two slots, $B_3$ in the last two, and the two strands of $B_2$ in the third and sixth slots.
The result is averaged over the three choices of the split argument, so that the form is symmetric in $(B_1, B_2, B_3)$.
Off the diagonal, this symmetrized form is the one meant whenever $\vertex{5 \cdot 5}$ is evaluated on $2$-string states, in particular in the reduced sextic vertex of \cref{p:sec:hs:sextic}.
On the diagonal this reading gives $\vertex{5 \cdot 5}(\Psi^2, \Psi^2, \Psi^2) = \vertex{5 \cdot 5}(\Psi^6)$.


Once the quintic vertex is included, the auxiliary field thus induces interactions at all orders.
A $\Sigma_2$ line collapses to a contact term when one of its ends is a source vertex, as in~\eqref{p:eq:hs4:telescoping}.
The line of $\vertex{5 \cdot 5}$ joins two quintic vertices and in general does not collapse.
The sextic order must therefore reproduce $\vertex{6}$ and cancel this induced term at the same time, which is the subject of \cref{p:sec:hs:sextic}.
\Cref{p:sec:worldsheet} gives the interpretation of this exchange in terms of moduli space.

This also means that this approach of using $\Sigma_2$ directly in the quintic term works only because SFT contains an infinite number of higher-order terms.
Indeed, any term induced at higher orders can be re-packaged with the higher-order elementary vertices.
In normal QFT, we would not be able to use the same technique, without also introducing an infinite tower of auxiliary fields.
Another way to deal with odd powers was suggested in~\cite{Rivasseau:2010:LoopVertexExpansion,Lionni:2019:IntermediateFieldRepresentation}.


\subparagraph{The five-point diagrams.}

The identity~\eqref{p:eq:hs5:identity} guarantees that every tree amplitude of the cubic action agrees with that of $S_{\le 5}$ through order $\gs^3$, under the assumptions on the exchange stated in \cref{p:sec:hs:quartic} and the standing assumptions of this subsection.
We show how the new diagram of the five-point amplitude reduces to the quintic vertex, since this is the first order at which the auxiliary field gives a new diagram.
Tree amplitudes with $n$ external legs scale as $\gs^{n - 2}$, so order $\gs^3$ is the five-point tree level.
The Feynman rules are those of \cref{p:sec:hs:quartic}, in the same convention: the propagator $\tfrac{12}{\gs^2} \sprod{2}{2}^{-1}$ and the source vertex $\tfrac{\gs^2}{3!} \vertex{4}(\Sigma_2, \Psi^2)$ of~\eqref{p:eq:hs4:feynman-rules}.
The quintic coupling~\eqref{p:eq:hs5:S5tilde} adds a vertex with two $\Sigma_2$ legs and one $\Psi$ leg:
\begin{equation}
  \label{p:eq:hs5:feynman-rule}
  \frac{2 \gs^3}{5!} \, \vertex{5}(\Sigma_2, \Sigma_2, \Psi).
\end{equation}
A $\Sigma_2$ line ends either on a source vertex, linear in $\Sigma_2$, or on the quintic vertex, quadratic in it.
\begin{figure}[htp]
  \centering
  \begin{tikzpicture}[baseline={(0, -0.1)}]
    \begin{feynman}
      \vertex[dot] (v) at (0, 0) {};
      \vertex (a) at (-1.2, 0.8) {$\Sigma_2$};
      \vertex (b) at (-1.2, -0.8) {$\Sigma_2$};
      \vertex (c) at (1.2, 0) {$\Psi$};
      \diagram* {
      (a) -- [double, double distance=1.5pt] (v) -- [double, double distance=1.5pt] (b),
      (v) -- [plain] (c),
      };
    \end{feynman}
  \end{tikzpicture}
  \qquad \qquad
  \begin{tikzpicture}[baseline={(0, -0.1)}]
    \begin{feynman}
      \vertex[dot] (s1) at (0, 0.9) {};
      \vertex[dot] (s2) at (0, -0.9) {};
      \vertex[dot] (q) at (1.6, 0) {};
      \vertex (p1) at (-1, 1.5) {$\Psi$};
      \vertex (p2) at (-1, 0.4) {$\Psi$};
      \vertex (p3) at (-1, -0.4) {$\Psi$};
      \vertex (p4) at (-1, -1.5) {$\Psi$};
      \vertex (p5) at (2.7, 0) {$\Psi$};
      \diagram* {
      (p1) -- [plain] (s1),
      (p2) -- [plain] (s1),
      (p3) -- [plain] (s2),
      (p4) -- [plain] (s2),
      (s1) -- [double, double distance=1.5pt, edge label={$\Sigma_2$}] (q),
      (s2) -- [double, double distance=1.5pt, edge label'={$\Sigma_2$}] (q),
      (q) -- [plain] (p5),
      };
    \end{feynman}
  \end{tikzpicture}
  \quad $\longrightarrow$ \quad
  \begin{tikzpicture}[baseline={(0, -0.1)}]
    \begin{feynman}
      \vertex[dot] (c) at (0, 0) {};
      \vertex (l1) at (90:1.1) {$\Psi$};
      \vertex (l2) at (162:1.1) {$\Psi$};
      \vertex (l3) at (234:1.1) {$\Psi$};
      \vertex (l4) at (306:1.1) {$\Psi$};
      \vertex (l5) at (18:1.1) {$\Psi$};
      \diagram* {
      (c) -- [plain] (l1),
      (c) -- [plain] (l2),
      (c) -- [plain] (l3),
      (c) -- [plain] (l4),
      (c) -- [plain] (l5),
      };
    \end{feynman}
  \end{tikzpicture}
  \caption{Left, the vertex~\eqref{p:eq:hs5:feynman-rule} of the quintic coupling, with two $\Sigma_2$ legs (double lines) and one $\Psi$ leg (single line); right, the five-point exchange diagram, two source vertices joined by $\Sigma_2$ lines to the quintic vertex, and its collapse to the contact vertex $\vertex{5}$.}
  \label{p:fig:hs5:five-point}
\end{figure}

The new diagram joins two source vertices by $\Sigma_2$ lines to the quintic vertex, whose $\Psi$ leg carries the fifth external state (\cref{p:fig:hs5:five-point}).
Both lines collapse at once, since the propagated source is the pair state: the exchange diagram evaluates $\vertex{5}(\Psi^2, \Psi^2, \Psi)$.
The collapse holds up to an element of $\ker \sprod{2}{2}$, on which the quintic vertex vanishes by~\eqref{p:eq:hs5:kernel-condition}.
A quintic vertex joined to a cubic vertex by its $\Psi$ leg, or to a second quintic vertex by a $\Sigma_2$ line, has six external legs and is of order $\gs^4$.
The second one is the diagram behind the induced term of~\eqref{p:eq:hs5:eliminate}.


\paragraph{Gauge transformation}

The transformation law~\eqref{p:eq:hs4:gauge-sigma} of $\Sigma_2$ at quartic order does not leave the action invariant at order $\gs^3$.

Since the solution~\eqref{p:eq:hs5:saddle} contains new terms which depend on $\Psi$, it transforms under~\eqref{p:eq:hs1:gauge-psi}, and the transformation of $\Sigma_2$ must be modified to compensate.
We generalize the pattern of the quartic order, where the transformation law~\eqref{p:eq:hs4:gauge-sigma} left the fluctuation invariant, and require the quadratic form $\pair[2]{\sigma_2^{(3)}}{\kin{2}{\Psi}{3}(\sigma_2^{(3)})}$ to be invariant.
We impose the invariance of this form because it is the only term of the identity~\eqref{p:eq:hs5:identity} which depends on $\Sigma_2$.
If it is invariant, the variation of the cubic action equals that of $S_{\le 5}$ up to the variation of the last term, which is of order $\gs^4$, the order of the truncation.
Its variation has two pieces, from the fluctuation and from the $\Psi$ inside the dressed operator.
By~\eqref{p:eq:hs5:dressed} and the symmetry of the dressed operator, we get:
\begin{equation}
  \label{p:eq:hs5:form-variation}
  \delta_\Lambda \pair[2]{\sigma_2^{(3)}}{\kin{2}{\Psi}{3}(\sigma_2^{(3)})}
    = 2 \, \pair[2]{\delta_\Lambda \sigma_2^{(3)}}{\kin{2}{\Psi}{3}(\sigma_2^{(3)})}
      - \frac{\gs}{5} \, \vertex{5}(\sigma_2^{(3)}, \sigma_2^{(3)}, \delta_\Lambda \Psi).
\end{equation}
This variation vanishes if the fluctuation transforms as:
\begin{equation}
  \label{p:eq:hs5:gauge-deviation}
  \delta_\Lambda \sigma_2^{(3)}
    = \frac{\gs}{10} \, \big( \kin{2}{\Psi}{3} \big)^{-1} \circ \sprod{3}{2}(\delta_\Lambda \Psi, \sigma_2^{(3)}).
\end{equation}
The fluctuation transforms homogeneously, by a $\Psi$-dependent linear rotation which vanishes on-shell.
This is the transformation law~\eqref{p:eq:gauge:law-quintic} of \cref{p:sec:gauge}, where its algebra is worked out.

Translating back to $\Sigma_2 = \Sigma_2^{(3)*} + \sigma_2^{(3)}$, the auxiliary field transforms as its own solution, plus the rotation of the fluctuation:
\begin{equation}
  \label{p:eq:hs5:gauge-sigma}
  \delta_\Lambda \Sigma_2
    = \delta_\Lambda \Sigma_2^{(3)*}
      + \frac{\gs}{10} \, \big( \kin{2}{\Psi}{3} \big)^{-1} \circ \sprod{3}{2}(\delta_\Lambda \Psi, \sigma_2^{(3)}),
\end{equation}
where $\delta_\Lambda \Sigma_2^{(3)*}$ is the variation of~\eqref{p:eq:hs5:saddle} induced by $\delta_\Lambda \Psi$.
The transformation law of the quartic order is the order-$\gs^0$ truncation of $\delta_\Lambda \Sigma_2 = \delta_\Lambda \Sigma_2^{(3)*}$.
Through order $\gs$, the transformation law reads:
\begin{equation}
  \label{p:eq:hs5:gauge-sigma-g1}
  \begin{aligned}
    \delta_\Lambda \Sigma_2
      = 2 \, \delta_\Lambda \Psi \otimes \Psi
        &
        + \frac{3 \gs}{5} \, \sprod{2}{2}^{-1} \circ \sprod{3}{2}(\delta_\Lambda \Psi, \Psi^2)
        \\ &
        + \frac{\gs}{10} \, \sprod{2}{2}^{-1} \circ \sprod{3}{2}(\delta_\Lambda \Psi, \Sigma_2 - \Psi^2)
        + O(\gs^2),
    \end{aligned}
\end{equation}
the middle term being the variation of $\tfrac{\gs}{5} \sprod{2}{2}^{-1} \circ \sprod{3}{2}(\Psi^3)$, cubic in $\Psi$.
The last term vanishes on-shell but cannot be dropped off-shell: it compensates the variation of the dressed operator.

With~\eqref{p:eq:hs5:gauge-sigma} the quadratic form is exactly invariant.
Varying the identity~\eqref{p:eq:hs5:identity} gives:
\begin{equation}
  \label{p:eq:hs5:delta-S-equal}
  \delta_\Lambda \caction[\le 5] - \delta_\Lambda S_{\le 5}
    = \frac{\gs^4}{600} \, \delta_\Lambda \vertex{5}\Big( \Psi^2, \Psi, \big( \kin{2}{\Psi}{3} \big)^{-1} \circ \sprod{3}{2}(\Psi^3) \Big)
    = O(\gs^4),
\end{equation}
identically, off-shell, for any variation $\delta_\Lambda \Psi$ and exactly in $\gs$.
Since $\delta_\Lambda S_{\le 5} = O(\gs^4)$, the cubic action is invariant up to terms of order $\gs^4$, the order of the truncation it should match.
We discuss the nature of the rotation term, the commutator of two gauge transformations and the Batalin--Vilkovisky structure of the cubic action in \cref{p:sec:gauge}.


\paragraph{Kernel condition and reality}

\subparagraph{Kernel condition.}

The auxiliary field lives on the quotient~\eqref{p:eq:hs4:quotient}, and the quintic coupling must be a function on it.
The two quartic terms of~\eqref{p:eq:hs5:F} are invariant under $\Sigma_2 \to \Sigma_2 + N_2$ for $N_2 \in \ker \sprod{2}{2}$, by~\eqref{p:eq:hs4:kernel}, while the quintic term shifts by
\begin{equation}
  \vertex{5}((\Sigma_2 + N_2)^2, \Psi) - \vertex{5}(\Sigma_2^2, \Psi)
    \to 2 \, \vertex{5}(N_2, \Sigma_2, \Psi) + \vertex{5}(N_2^2, \Psi).
\end{equation}
For this to vanish for all $\Sigma_2$, the term linear in $\Sigma_2$ must vanish, which then kills the quadratic one.
The cubic action therefore descends to the quotient if and only if the quintic vertex vanishes on the kernel:
\begin{equation}
  \label{p:eq:hs5:kernel-condition}
  \forall N_2 \in \ker \sprod{2}{2}, \;
  \forall A \in \Hn{2}, \;
  \forall B \in \Hs:
  \quad
  \vertex{5}(N_2, A, B)
    = 0,
\end{equation}
that is, $\sprod{3}{2}(A, B)$ pairs to zero with every kernel element.
In infinite dimension, the states which pair to zero with the kernel of $\sprod{2}{2}$ form only the closure of its image: the condition thus places the values of $\sprod{3}{2}$ in the closure of the image of $\sprod{2}{2}$, and the first assumption below, that they lie in the image itself, is what makes $\sprod{2}{2}^{-1}$ defined on them.
It also says that $\sprod{3}{2}(\Psi, \cdot)$ does not depend on the representative of its argument.
The condition thus makes $\sprod{3}{2}(\Psi, \cdot)$ descend to the quotient, and with it the dressed operator~\eqref{p:eq:hs5:dressed-factor} and its inverse~\eqref{p:eq:hs5:neumann}.
Since $\vertex{5}$ is symmetric, the condition is equivalent by polarization to its restriction to the diagonal, $\vertex{5}(N_2, \Psi^2, \Psi) = 0$ for all $\Psi$, which is the only evaluation of the quintic vertex in the matching at order $\gs^3$.

We have no proof that the quintic vertex vanishes on a null vector of $\sprod{2}{2}$, and nothing imposes it.
  At the same time, no null vector of $\sprod{2}{2}$ is known above level $0$ (at level $0$ there is none at any real momentum for the sections with $\beta > \sqrt{2}$, \cref{p:sec:hs:quartic}), and none is forced by the mechanisms we have checked at the levels and momenta where the couplings are known: the dilaton theorem with the ghost-number counting behind the zeros of the cubic vertex, BRST-exact strands, the universality of the zero-momentum couplings, and the worldsheet twist.
  If $\sprod{2}{2}$ has an accidental zero at some real momentum, with a single null vector $N_2$, the condition at that momentum is that $N_2$ pairs to zero with the quintic source $\sprod{3}{2}(\Psi^2, \Psi)$, one number for each $\Psi$, a property of $\vertex{5}$ at that momentum which nothing enforces and which we expect to fail in general; the matching is then restored by the multiplier described next.

When it fails, the quotient is not available and one must choose a complement of the kernel.
The eliminated action then differs from $S_{\le 5}$ at order $\gs^3$ by a five-point term which depends on the complement and cannot vanish for all of them.
The kernel directions can then be handled by a multiplier rather than by inversion, which restores the matching at this order and keeps the action cubic; this is taken up in \cref{p:sec:all-orders}.

Hence, unless said otherwise, we make two assumptions from this order on: 1) the states on which $\sprod{2}{2}^{-1}$ acts, here the values of $\sprod{3}{2}$, lie in the image of $\sprod{2}{2}$ and not only in its closure, which is all the kernel condition gives, so that the inverse exists on them level by level (\cref{p:sec:hs:quartic}); 2) the quintic vertex vanishes on the kernel of the quartic vertex,~\eqref{p:eq:hs5:kernel-condition} (``kernel condition'').
The sextic vertex is subject to an analogous condition at the next order.


\subparagraph{Reality.}

The quintic order is the first where the solution and the gauge law carry an inverse, so that reality is trickier to prove.

We use the reality of the vertices~\eqref{p:eq:hs4:vertex-real} for $n = 4$ and $n = 5$, and the condition~\eqref{p:eq:hs5:kernel-condition}.
  The kernel of $\sprod{2}{2}$ is stable under the star conjugation.
  Indeed, if $\vertex{4}(A, N_2) = 0$ for all $A$, then $\vertex{4}(A, N_2^\star) = \vertex{4}(A^\star, N_2)^* = 0$ for all $A$, since $A^\star$ ranges over all $2$-string states as $A$ does.
  Hence $\star$ descends to the quotient, and the condition~\eqref{p:eq:hs4:reality}, $\Sigma_2^\star = \Sigma_2$, makes sense there.

The two maps from which the solution is built, $\sprod{2}{2}^{-1} \circ \sprod{3}{2}$ and $(\kin{2}{\Psi}{3})^{-1} \circ \sprod{2}{2}$, commute with $\star$ for real $\Psi$.
This follows by conjugating both sides of the two relations which characterize them:
\begin{equation}
\vertex{4}\big( A, \sprod{2}{2}^{-1} \circ \sprod{3}{2}(C) \big)
  = \vertex{5}(A, C),
\qquad
\pair[2]{A}{\kin{2}{\Psi}{3}(B')}
  = \vertex{4}(A, B),
\end{equation}
with $B' := (\kin{2}{\Psi}{3})^{-1} \circ \sprod{2}{2}(B)$.
The conjugation is applied only to states whose strands have the ghost number of the string field, in $\Hs$, $\Hn{2}$ and $\Hn{3}$, where it acts strand by strand.
It is never applied to the outputs of $\sprod{2}{2}$ and $\sprod{3}{2}$, whose strands have ghost number three.
All coefficients are real rationals times $\gs$.

With $\Psi^\star = - \Psi$, the state $\Psi \otimes \Psi$ is star-even.
For $\gs$ real, the solution~\eqref{p:eq:hs5:saddle} therefore satisfies $\Sigma_2^\star = \Sigma_2$, to all orders of the formal expansion in $\gs$ and, within the quintic truncation, wherever the dressed operator is invertible on the quotient.
For a real variation $\delta_\Lambda \Psi$, each term of the transformation law~\eqref{p:eq:hs5:gauge-sigma} is real for the same reason, and the transformation law maps the real locus of $(\Psi, \Sigma_2)$ to itself.
A gauge parameter satisfying the reality condition $\Lambda^\star = - \Lambda$ of \cref{p:app:conventions} produces such a real variation, the vertices being real.
